\documentclass[../main.tex]{subfiles}

\begin{document}

\section{Validazione del modulo} \label{sec:validazione}
  Per validare la correttezza del modulo realizzato, sono stati creati 19
  spill simulando $\sim$$10^{14}$ POT (\textit{Proton On Target}) tramite la
  pipeline di simulazione riportata in figura~\ref{fig:nd-sim} che sono poi
  stati copiati in un CAF da fake-reco come mostrato nello schema in
  figura~\ref{fig:fake-reco-v0}.
  Graficando le componenti y e z dei vertici nel volume attivo di SAND contenuti nei campi della
  ricostruzione si è ottenuto il grafico
  riportato in figura~\ref{fig:plt-hit}.

  \begin{figure}[ht]
    \centering
    \includegraphics[width=.95\textwidth]{fake-reco/plot-hit}
    \caption{
      Grafico dei vertici di interazione in SAND di 19 spill generati da
      $\sim$$10^{14}$ POT (\textit{Proton On Target}) contenuti nei campi
      della ricostruzione in un file CAF\@.
      I dati sono stati generati dalla catena di simulazione del
      \textit{Near Detector} e copiati nel CAF con fake-reco.
    }
    \label{fig:plt-hit}
  \end{figure}

\end{document}